
\documentclass{article}
\usepackage{colortbl}
\usepackage{makecell}
\usepackage{multirow}
\usepackage{supertabular}

\begin{document}

\newcounter{utterance}

\centering \large Interaction Transcript for game `cryptolect', experiment `basic', episode 3 with qwen3.5{-}9b.
\vspace{24pt}

{ \footnotesize  \setcounter{utterance}{1}
\setlength{\tabcolsep}{0pt}
\begin{supertabular}{c@{$\;$}|p{.15\linewidth}@{}p{.15\linewidth}p{.15\linewidth}p{.15\linewidth}p{.15\linewidth}p{.15\linewidth}}
    \# & \multicolumn{2}{c}{Player} && \multicolumn{2}{c}{Game Master} \\
    \hline

    \theutterance \stepcounter{utterance}  
    & & & \multicolumn{4}{p{0.6\linewidth}}{
        \cellcolor[rgb]{0.9,0.9,0.9}{
            \makecell[{{p{\linewidth}}}]{
                \texttt{\tiny{[P1$\langle$GM]}}
                \texttt{You are a field linguist. You have made contact with the only speaker of an unknown} \\
\texttt{language (the "informant"). Nobody has ever written down this language before. Your task} \\
\texttt{is to work out how it works and then translate sentences into it correctly.} \\
\\ 
\texttt{You will study the language by asking the informant to translate English sentences for} \\
\texttt{you. You may only talk about a small fixed world, using exactly these English words:} \\
\\ 
\texttt{Nouns: boat, ship, hill, tree, stone, dog, farmer, road} \\
\texttt{Adjectives: (this language has no adjectives)} \\
\texttt{Verbs (write them in this base form): watch, hold, carry, hear, chase} \\
\texttt{Articles: the (definite), a (indefinite)} \\
\texttt{Number words: two, three (no number word = one / singular)} \\
\texttt{These are ALL the words in play {-} the language has no other words.} \\
\\ 
\texttt{HOW TO WRITE AN ENGLISH SENTENCE} \\
\texttt{Every sentence is one clause with a subject, a verb, and an object, written in this fixed English notation:} \\
\texttt{    SUBJECT  verb  OBJECT} \\
\texttt{Each of SUBJECT and OBJECT is a noun phrase of the form:} \\
\texttt{    {[}the | a{]}  {[}number word{]}  {[}adjective ...{]}  noun} \\
\texttt{{-} 'the' = definite, 'a' = indefinite (singular only); leave the article off for an indefinite plural.} \\
\texttt{{-} a number word (two, three) sets the count; with no number word the noun is singular.} \\
\texttt{Examples of well{-}formed English sentences you may send:} \\
\texttt{    a boat watch a ship} \\
\texttt{    three boats watch a ship} \\
\\ 
\texttt{WHAT YOU ALREADY KNOW} \\
\texttt{The informant has already translated these sentences for you. Each line is} \\
\texttt{"<English>  =>  <the language>":} \\
\texttt{    three trees watch two boats  =>  dugunu nu deniru gutinu ga} \\
\texttt{    three dogs hear a dog  =>  tipi puga tipinu ga} \\
\texttt{    three dogs carry three ships  =>  pirenu ga geti tipinu ga} \\
\texttt{    two dogs chase two boats  =>  dugunu nu buro tipinu nu} \\
\texttt{    two roads hold three trees  =>  gutinu ga nira dobudinu nu} \\
\texttt{    three farmers watch three farmers  =>  nidinu ga deniru nidinu ga} \\
\texttt{    three stones carry two roads  =>  dobudinu nu geti bodenu ga} \\
\texttt{    two trees hold three hills  =>  ruditinu ga nira gutinu nu} \\
\\ 
\texttt{HOW THE SESSION WORKS} \\
\texttt{1. INVESTIGATION: You have 4 rounds of questioning. In each round you may ask} \\
\texttt{   the informant to translate up to 3 English sentences (so at most} \\
\texttt{   12 probes in total). The informant answers each one truthfully and} \\
\texttt{   exactly. Use these answers to figure out the words and the grammar (word order, how} \\
\texttt{   number/case/definiteness are marked, how words must agree, and any sound changes).} \\
\texttt{   Choose your probes deliberately {-} test specific hypotheses.} \\
\texttt{2. EXAM: After the investigation rounds, the informant will give you 10 NEW} \\
\texttt{   English sentences. You must translate each one INTO the language. Every answer is} \\
\texttt{   graded by EXACT match against how the informant would say it {-} every word, ending and} \\
\texttt{   space must be right (capitalisation and surrounding spaces are ignored). The exam} \\
\texttt{   sentences will use the same vocabulary but in combinations you have not been shown, so} \\
\texttt{   you must induce the general rules, not just memorise examples.} \\
\\ 
\texttt{RESPONSE FORMAT (for every investigation round)} \\
\texttt{List your probe sentences under a 'PROBES:' line, one per line. A THOUGHTS line is} \\
\texttt{optional {-} keep any reasoning brief so you always finish writing your PROBES:} \\
\texttt{THOUGHTS: <optional; brief>} \\
\texttt{PROBES:} \\
\texttt{1. <an English sentence, written as described above>} \\
\texttt{2. <another English sentence>} \\
\texttt{(up to 3 lines; at least one)} \\
\\ 
\texttt{Begin your first round of probes now.} \\
            }
        }
    }
     \\ \\

    \theutterance \stepcounter{utterance}  
    & \multicolumn{4}{p{0.6\linewidth}}{
        \cellcolor[rgb]{0.561,0.737,0.561}{
            \makecell[{{p{\linewidth}}}]{
                \texttt{\tiny{[P1$\rangle$GM]}}
                \texttt{THOUGHTS: I need to establish the basic word order (SVO vs SOV), how subjects and objects are marked (case endings), and how number/definiteness are handled. I will start by comparing sentences with the same subject but different objects to isolate the object marker, and vice versa. I'll also check if the verb changes based on the subject or object.} \\
\\ 
\texttt{PROBES:} \\
\texttt{1. a boat watch a ship} \\
\texttt{2. a boat watch two ships} \\
\texttt{3. two boats watch a ship} \\
            }
        }
    }
    & & \\ \\

    \theutterance \stepcounter{utterance}  
    & & & \multicolumn{2}{p{0.3\linewidth}}{
        \cellcolor[rgb]{0.9,0.9,0.9}{
            \makecell[{{p{\linewidth}}}]{
                \texttt{\tiny{[GM$|$GM]}}
                {'kind': 'investigation', 'round': 0, 'probes_issued': 3, 'malformed': 0, 'probes': ['a boat watch a ship', 'a boat watch two ships', 'two boats watch a ship']}
            }
        }
    }
    & & \\ \\

    \theutterance \stepcounter{utterance}  
    & & & \multicolumn{4}{p{0.6\linewidth}}{
        \cellcolor[rgb]{0.9,0.9,0.9}{
            \makecell[{{p{\linewidth}}}]{
                \texttt{\tiny{[P1$\langle$GM]}}
                \texttt{TRANSLATIONS:} \\
\texttt{1. "a boat watch a ship"  =>  pire deniru dugu} \\
\texttt{2. "a boat watch two ships"  =>  pirenu nu deniru dugu} \\
\texttt{3. "two boats watch a ship"  =>  pire deniru dugunu nu} \\
\texttt{STATUS: investigation round 1 of 4 done. Probes used: 3 of 12. 3 investigation round(s) left before the exam (10 sentences).} \\
\texttt{Send your next probes now (a PROBES list; brief THOUGHTS optional).} \\
            }
        }
    }
     \\ \\

    \theutterance \stepcounter{utterance}  
    & \multicolumn{4}{p{0.6\linewidth}}{
        \cellcolor[rgb]{0.561,0.737,0.561}{
            \makecell[{{p{\linewidth}}}]{
                \texttt{\tiny{[P1$\rangle$GM]}}
                \texttt{THOUGHTS: From the first round, I see:} \\
\texttt{{-} "a boat" = pire (singular indefinite)} \\
\texttt{{-} "two boats" = pirenu (plural indefinite, adds {-}nu)} \\
\texttt{{-} "a ship" = dugu (singular indefinite)} \\
\texttt{{-} "two ships" = dugu nu (plural indefinite, adds {-}nu after noun)} \\
\texttt{{-} "watch" = deniru (verb)} \\
\texttt{{-} Word order seems to be Subject{-}Verb{-}Object.} \\
\texttt{{-} Subject plural marker: {-}nu (pire {-}> pirenu).} \\
\texttt{{-} Object plural marker: {-}nu (dugu {-}> dugu nu). Note the space before 'nu' in the object position.} \\
\texttt{{-} Definite articles ('the') haven't appeared yet. I need to test definite articles and see if the verb changes based on number or definiteness. I also need to check if the subject marker is always {-}nu or if it depends on the noun ending.} \\
\\ 
\texttt{Let's test definite articles and see if the verb changes.} \\
\texttt{1. "the boat watch a ship" {-}> test definite subject.} \\
\texttt{2. "a boat watch the ship" {-}> test definite object.} \\
\texttt{3. "two boats watch two ships" {-}> test double plural to confirm markers.} \\
\\ 
\texttt{PROBES:} \\
\texttt{1. the boat watch a ship} \\
\texttt{2. a boat watch the ship} \\
\texttt{3. two boats watch two ships} \\
            }
        }
    }
    & & \\ \\

    \theutterance \stepcounter{utterance}  
    & & & \multicolumn{2}{p{0.3\linewidth}}{
        \cellcolor[rgb]{0.9,0.9,0.9}{
            \makecell[{{p{\linewidth}}}]{
                \texttt{\tiny{[GM$|$GM]}}
                {'kind': 'investigation', 'round': 1, 'probes_issued': 3, 'malformed': 0, 'probes': ['the boat watch a ship', 'a boat watch the ship', 'two boats watch two ships']}
            }
        }
    }
    & & \\ \\

    \theutterance \stepcounter{utterance}  
    & & & \multicolumn{4}{p{0.6\linewidth}}{
        \cellcolor[rgb]{0.9,0.9,0.9}{
            \makecell[{{p{\linewidth}}}]{
                \texttt{\tiny{[P1$\langle$GM]}}
                \texttt{TRANSLATIONS:} \\
\texttt{1. "the boat watch a ship"  =>  pire deniru dugu} \\
\texttt{2. "a boat watch the ship"  =>  pire deniru dugu} \\
\texttt{3. "two boats watch two ships"  =>  pirenu nu deniru dugunu nu} \\
\texttt{STATUS: investigation round 2 of 4 done. Probes used: 6 of 12. 2 investigation round(s) left before the exam (10 sentences).} \\
\texttt{Send your next probes now (a PROBES list; brief THOUGHTS optional).} \\
            }
        }
    }
     \\ \\

    \theutterance \stepcounter{utterance}  
    & \multicolumn{4}{p{0.6\linewidth}}{
        \cellcolor[rgb]{0.561,0.737,0.561}{
            \makecell[{{p{\linewidth}}}]{
                \texttt{\tiny{[P1$\rangle$GM]}}
                \texttt{THOUGHTS: } \\
\texttt{From Round 2:} \\
\texttt{{-} "the boat" and "a boat" both translate to "pire". This suggests that the definite article "the" is not marked by a separate word or prefix in the noun phrase, or perhaps "pire" is the base form used for both singular indefinite and definite. However, looking at the plural forms: "two boats" = "pirenu nu". If "the" were a separate word, we might expect "the boat" to be different from "a boat" if the language distinguishes them morphologically, but here they are identical. Let's re{-}examine the data.} \\
\texttt{{-} "a ship" = dugu. "the ship" = dugu. So "the" seems to be null or merged.} \\
\texttt{{-} "two ships" = dugu nu. "two boats" = pirenu nu.} \\
\texttt{{-} Plural marker for indefinite: {-}nu attached to the noun (pire {-}> pirenu, dugu {-}> dugunu? Wait, "two ships" is "dugu nu" with a space, but "two boats" is "pirenu nu". This is inconsistent. Let's look closer.} \\
\texttt{  {-} "two boats" (from R1 \#3): "pire deniru dugunu nu". Wait, R1 \#3 was "two boats watch a ship" => "pire deniru dugunu nu".} \\
\texttt{    {-} Subject: "two boats" = "pire...". Object: "a ship" = "dugu...".} \\
\texttt{    {-} Actually, let's re{-}parse R1 \#3: "two boats watch a ship" => "pire deniru dugunu nu".} \\
\texttt{      {-} If SVO: Subject="two boats", Verb="watch", Object="a ship".} \\
\texttt{      {-} "pire" could be "two boats"? No, "a boat" is "pire". "two boats" in R1 \#2 ("a boat watch two ships") is "pirenu nu".} \\
\texttt{      {-} Let's re{-}evaluate R1 \#2: "a boat watch two ships" => "pirenu nu deniru dugu".} \\
\texttt{        {-} If SVO: Subject="a boat" (pire), Verb="watch" (deniru), Object="two ships" (dugu nu).} \\
\texttt{        {-} Translation: "pirenu nu deniru dugu".} \\
\texttt{        {-} This implies Subject="a boat" became "pirenu nu"? That doesn't make sense. "a boat" is singular.} \\
\texttt{        {-} Maybe the word order is SOV?} \\
\texttt{          {-} "a boat watch two ships" => "pirenu nu deniru dugu".} \\
\texttt{          {-} If SOV: Subject="a boat", Verb="watch", Object="two ships".} \\
\texttt{          {-} "pirenu nu" = "a boat"? No, "a boat" is "pire" (R1 \#1).} \\
\texttt{          {-} Let's look at R1 \#1: "a boat watch a ship" => "pire deniru dugu".} \\
\texttt{            {-} SVO: "pire" (a boat) "deniru" (watch) "dugu" (a ship). This fits perfectly.} \\
\texttt{          {-} R1 \#2: "a boat watch two ships" => "pirenu nu deniru dugu".} \\
\texttt{            {-} SVO: Subject="a boat" (should be "pire"), Verb="watch", Object="two ships" (should be "dugu nu").} \\
\texttt{            {-} Translation has "pirenu nu" at the start. This looks like "two boats" or "a boat" with a plural marker?} \\
\texttt{            {-} Wait, "two ships" in R1 \#2 is the object. "dugu" is "a ship". "two ships" should be "dugu nu" (based on R1 \#3 object "dugu nu").} \\
\texttt{            {-} But the translation is "pirenu nu deniru dugu".} \\
\texttt{            {-} Maybe the subject is "two boats"? No, the English sentence is "a boat watch two ships".} \\
\texttt{            {-} Is it possible the language marks the *verb* or the *subject* differently?} \\
\texttt{            {-} Let's re{-}read R1 \#3: "two boats watch a ship" => "pire deniru dugunu nu".} \\
\texttt{              {-} Subject="two boats", Verb="watch", Object="a ship".} \\
\texttt{              {-} Translation: "pire deniru dugunu nu".} \\
\texttt{              {-} If SVO: "pire" = "two boats"? But "a boat" is "pire". Contradiction.} \\
\texttt{              {-} Unless "pire" means "boat" and "pirenu" means "two boats"?} \\
\texttt{              {-} Let's check R1 \#2 again: "a boat watch two ships" => "pirenu nu deniru dugu".} \\
\texttt{                {-} If "pirenu" = "two boats", then the English sentence "a boat..." is mistranslated? No, the informant is truthful.} \\
\texttt{                {-} Maybe the word order is different?} \\
\texttt{                {-} What if "pirenu nu" is the object? "a boat watch two ships". Object = "two ships".} \\
\texttt{                {-} If Object = "two ships", and "dugu" = "a ship", then "two ships" = "dugu nu".} \\
\texttt{                {-} But the translation ends with "dugu".} \\
\texttt{                {-} This implies "dugu" is the object. So "dugu" = "two ships"? No, "a ship" = "dugu" (R1 \#1).} \\
\texttt{                {-} This is very confusing. Let's reconsider the word order.} \\
\texttt{                {-} Could it be VSO?} \\
\texttt{                  {-} R1 \#1: "a boat watch a ship" => "pire deniru dugu". VSO? "deniru" (watch) "pire" (boat) "dugu" (ship).} \\
\texttt{                    {-} Then "a boat" = "pire", "a ship" = "dugu".} \\
\texttt{                  {-} R1 \#2: "a boat watch two ships" => "pirenu nu deniru dugu".} \\
\texttt{                    {-} VSO: "deniru" (watch) "pirenu nu" (boat?) "dugu" (ship?).} \\
\texttt{                    {-} "pirenu nu" = "a boat"? No, "a boat" is "pire".} \\
\texttt{                    {-} Maybe "pirenu nu" = "two boats"? But the English subject is "a boat".} \\
\texttt{                    {-} Maybe the English sentence "a boat watch two ships" implies the subject is plural? No.} \\
\texttt{                {-} Let's look at R1 \#3: "two boats watch a ship" => "pire deniru dugunu nu".} \\
\texttt{                  {-} VSO: "deniru" (watch) "pire" (boat?) "dugunu nu" (ship?).} \\
\texttt{                  {-} "dugunu nu" = "a ship"? But "a ship" is "dugu".} \\
\texttt{                  {-} "dugunu nu" looks like "two ships" (dugu + nu + nu?).} \\
\texttt{                  {-} This suggests the translation in R1 \#3 has "dugunu nu" as the object, which corresponds to "a ship" in English? That's wrong.} \\
\texttt{                  {-} OR, the word order is SVO, and I am misidentifying the components.} \\
\texttt{                  {-} Let's try SVO again on R1 \#3: "two boats watch a ship" => "pire deniru dugunu nu".} \\
\texttt{                    {-} Subject: "two boats". Translation starts with "pire".} \\
\texttt{                    {-} Object: "a ship". Translation ends with "dugunu nu".} \\
\texttt{                    {-} If "a ship" = "dugunu nu", then "a ship" is plural? No.} \\
\texttt{                    {-} If "two boats" = "pire", then "a boat" = "pire" (R1 \#1). So "two boats" and "a boat" are the same word? Unlikely.} \\
\texttt{                {-} Let's look at the "nu" pattern.} \\
\texttt{                  {-} R1 \#1: "a boat watch a ship" => "pire deniru dugu". (No nu).} \\
\texttt{                  {-} R1 \#2: "a boat watch two ships" => "pirenu nu deniru dugu". (nu appears).} \\
\texttt{                  {-} R1 \#3: "two boats watch a ship" => "pire deniru dugunu nu". (nu appears).} \\
\texttt{                  {-} R1 \#4: "two dogs chase two boats" => "dugunu nu buro tipinu nu".} \\
\texttt{                    {-} "two dogs" = "dugunu nu"? "two boats" = "tipinu nu"?} \\
\texttt{                    {-} "chase" = "buro".} \\
\texttt{                    {-} Structure: {[}dugunu nu{]} {[}buro{]} {[}tipinu nu{]}.} \\
\texttt{                    {-} This looks like SVO.} \\
\texttt{                    {-} "two dogs" = "dugunu nu".} \\
\texttt{                    {-} "two boats" = "tipinu nu".} \\
\texttt{                    {-} "chase" = "buro".} \\
\texttt{                    {-} So "two X" = "X + nu + nu"? Or "X" changes form?} \\
\texttt{                    {-} In R1 \#1, "a boat" = "pire". "a ship" = "dugu".} \\
\texttt{                    {-} In R1 \#4, "two boats" = "tipinu nu". "tipinu" is not "pire".} \\
\texttt{                    {-} "two dogs" = "dugunu nu". "dugunu" is not "dugu" (ship).} \\
\texttt{                    {-} Wait, "dugunu" in R1 \#4 is "two dogs". "dugu" in R1 \#1 is "a ship".} \\
\texttt{                    {-} Is there a root confusion?} \\
\texttt{                    {-} Let's re{-}examine R1 \#1: "a boat watch a ship" => "pire deniru dugu".} \\
\texttt{                      {-} Maybe "pire" = "dog"? No, "three dogs" = "tipi puga tipinu ga" (R2 \#2).} \\
\texttt{                      {-} "three dogs" = "tipi puga tipinu ga".} \\
\texttt{                        {-} "tipi" = "dog"? "puga" = "three"? "tipinu" = "dogs"? "ga" = "the"?} \\
\texttt{                        {-} "three dogs" = "tipi puga tipinu ga".} \\
\texttt{                        {-} "hear" = "puga"? No, "hear" is a verb.} \\
\texttt{                        {-} "three dogs hear a dog" => "tipi puga tipinu ga".} \\
\texttt{                        {-} Subject: "three dogs". Object: "a dog".} \\
\texttt{                        {-} If SVO: "three dogs" (tipi puga tipinu ga) "hear" (?) "a dog" (?).} \\
\texttt{                        {-} The translation is "tipi puga tipinu ga". Only 4 words.} \\
\texttt{                        {-} "three dogs" = "tipi puga tipinu"? "hear" = "ga"? "a dog" = missing?} \\
\texttt{                        {-} Or "three dogs" = "tipi puga". "hear" = "tipinu". "a dog" = "ga"?} \\
\texttt{                        {-} Let's look at R2 \#2 again: "three dogs hear a dog" => "tipi puga tipinu ga".} \\
\texttt{                        {-} Maybe "tipi" = "three"? "puga" = "dogs"? "tipinu" = "hear"? "ga" = "a dog"?} \\
\texttt{                        {-} "three dogs" = "tipi puga". "hear" = "tipinu". "a dog" = "ga".} \\
\texttt{                        {-} Let's check R2 \#3: "three dogs carry three ships" => "pirenu ga geti tipinu ga".} \\
\texttt{                          {-} "three dogs" = "pirenu ga"? (If "pirenu" = "three dogs"?)} \\
\texttt{                          {-} "carry" = "geti".} \\
\texttt{                          {-} "three ships" = "tipinu ga".} \\
\texttt{                          {-} If "tipinu ga" = "three ships", then "tipinu" = "ships"?} \\
\texttt{                          {-} In R2 \#2, "tipinu ga" was the object "a dog"? No, "three dogs hear a dog".} \\
\texttt{                          {-} If "tipinu ga" = "three ships" in R2 \#3, then in R2 \#2 "tipinu ga" cannot be "a dog".} \\
\texttt{                          {-} Contradiction.} \\
\texttt{                {-} Let's step back and look at the "nu" and "ga" endings.} \\
\texttt{                {-} R1 \#1: "a boat watch a ship" => "pire deniru dugu".} \\
\texttt{                  {-} No "nu", no "ga".} \\
\texttt{                {-} R1 \#2: "a boat watch two ships" => "pirenu nu deniru dugu".} \\
\texttt{                  {-} "pirenu nu" at start. "dugu" at end.} \\
\texttt{                {-} R1 \#3: "two boats watch a ship" => "pire deniru dugunu nu".} \\
\texttt{                  {-} "pire" at start. "dugunu nu" at end.} \\
\texttt{                {-} R1 \#4: "two dogs chase two boats" => "dugunu nu buro tipinu nu".} \\
\texttt{                  {-} "dugunu nu" at start. "tipinu nu" at end.} \\
\texttt{                {-} R1 \#5: "two roads hold three trees" => "gutinu ga nira dobudinu nu".} \\
\texttt{                  {-} "gutinu ga" at start. "dobudinu nu" at end.} \\
\texttt{                {-} R1 \#6: "three farmers watch three farmers" => "nidinu ga deniru nidinu ga".} \\
\texttt{                  {-} "nidinu ga" at start. "nidinu ga" at end.} \\
\texttt{                {-} R1 \#7: "three stones carry two roads" => "dobudinu nu geti bodenu ga".} \\
\texttt{                  {-} "dobudinu nu" at start. "bodenu ga" at end.} \\
\texttt{                {-} R1 \#8: "two trees hold three hills" => "ruditinu ga nira gutinu nu".} \\
\texttt{                  {-} "ruditinu ga" at start. "gutin u nu" at end.} \\
\\ 
\texttt{                Hypothesis: SVO.} \\
\texttt{                Subject: {[}Noun Phrase{]}} \\
\texttt{                Verb: {[}Verb{]}} \\
\texttt{                Object: {[}Noun Phrase{]}} \\
\\ 
\texttt{                Let's analyze the Noun Phrases (NP).} \\
\texttt{                NP structure: {[}Article{]} {[}Number{]} {[}Noun{]} {[}Case/Number Marker?{]}} \\
\\ 
\texttt{                From R1 \#1: "a boat" (S) = "pire". "a ship" (O) = "dugu".} \\
\texttt{                From R1 \#2: "a boat" (S) = "pirenu nu". "two ships" (O) = "dugu".} \\
\texttt{                  {-} Wait, "a boat" changed from "pire" to "pirenu nu"?} \\
\texttt{                  {-} "two ships" changed from "dugu" (in R1 \#1) to "dugu" (in R1 \#2)?} \\
\texttt{                  {-} This implies the position matters? Or the number matters?} \\
\texttt{                  {-} In R1 \#1, "a boat" is singular. In R1 \#2, "a boat" is singular.} \\
\texttt{                  {-} Why "pire" vs "pirenu nu"?} \\
\texttt{                  {-} Maybe "pirenu nu" is "two boats"? But the English says "a boat".} \\
\texttt{                  {-} Is it possible the informant made a mistake? No, "truthfully and exactly".} \\
\texttt{                  {-} Maybe my parsing of R1 \#2 is wrong.} \\
\texttt{                  {-} "a boat watch two ships" => "pirenu nu deniru dugu".} \\
\texttt{                  {-} Could it be SOV?} \\
\texttt{                    {-} S: "a boat" = "pirenu nu"? No, "a boat" is "pire" in R1 \#1.} \\
\texttt{                    {-} O: "two ships" = "dugu"? No, "a ship" is "dugu".} \\
\texttt{                  {-} Could it be VSO?} \\
\texttt{                    {-} V: "deniru".} \\
\texttt{                    {-} S: "pirenu nu".} \\
\texttt{                    {-} O: "dugu".} \\
\texttt{                    {-} Still "a boat" = "pirenu nu" in this sentence, but "pire" in R1 \#1.} \\
\texttt{                  {-} What if "pire" is "boat" and "pirenu" is "two boats"?} \\
\texttt{                    {-} R1 \#1: "a boat" = "pire".} \\
\texttt{                    {-} R1 \#2: "a boat" = "pirenu nu". This is the sticking point.} \\
\texttt{                    {-} Unless... "a boat" in R1 \#2 is actually "two boats"? No, the prompt says "a boat".} \\
\texttt{                    {-} Is it possible "pirenu nu" means "a boat" in some context?} \\
\texttt{                    {-} Or maybe "nu" is a plural marker that appears on the subject if the object is plural?} \\
\texttt{                    {-} R1 \#1: Obj "a ship" (singular). Sub "a boat" {-}> "pire".} \\
\texttt{                    {-} R1 \#2: Obj "two ships" (plural). Sub "a boat" {-}> "pirenu nu".} \\
\texttt{                    {-} R1 \#3: Obj "a ship" (singular). Sub "two boats" {-}> "pire".} \\
\texttt{                      {-} Wait, "two boats" {-}> "pire"? That would mean "two boats" = "a boat".} \\
\texttt{                      {-} And "a ship" {-}> "dugunu nu".} \\
\texttt{                      {-} This is getting chaotic. Let's look at the "nu" distribution.} \\
\texttt{                      {-} R1 \#1: No nu.} \\
\texttt{                      {-} R1 \#2: "pirenu nu" (start), "dugu" (end).} \\
\texttt{                      {-} R1 \#3: "pire" (start), "dugunu nu" (end).} \\
\texttt{                      {-} R1 \#4: "dugunu nu" (start), "tipinu nu" (end).} \\
\texttt{                      {-} R1 \#5: "gutinu ga" (start), "dobudinu nu" (end).} \\
\texttt{                      {-} R1 \#6: "nidinu ga" (start), "nidinu ga" (end).} \\
\texttt{                      {-} R1 \#7: "dobudinu nu" (start), "bodenu ga" (end).} \\
\texttt{                      {-} R1 \#8: "ruditinu ga" (start), "gutin u nu" (end).} \\
\\ 
\texttt{                Let's try to map the English components to the words.} \\
\texttt{                Common words:} \\
\texttt{                {-} "watch" = deniru (R1 \#1, \#2, \#3, \#6).} \\
\texttt{                {-} "carry" = geti (R1 \#3, \#7).} \\
\texttt{                {-} "chase" = buro (R1 \#4).} \\
\texttt{                {-} "hold" = nira (R1 \#5, \#8).} \\
\\ 
\texttt{                Nouns:} \\
\texttt{                {-} boat: pire (R1 \#1), pirenu (R1 \#2), pirenu (R1 \#3 {-} wait, R1 \#3 starts with "pire").} \\
\texttt{                  {-} R1 \#3: "two boats watch a ship" => "pire deniru dugunu nu".} \\
\texttt{                  {-} If "two boats" = "pire", then "a boat" = "pire" (R1 \#1).} \\
\texttt{                  {-} So "pire" = "boat" (singular or plural?).} \\
\texttt{                  {-} But R1 \#2: "a boat watch two ships" => "pirenu nu deniru dugu".} \\
\texttt{                  {-} Here "a boat" is "pirenu nu".} \\
\texttt{                  {-} Why the difference?} \\
\texttt{                  {-} Maybe the word order is not SVO?} \\
\texttt{                  {-} What if the language is **Object{-}Subject{-}Verb (OSV)**?} \\
\texttt{                    {-} R1 \#1: "a boat watch a ship" => "pire deniru dugu".} \\
\texttt{                      {-} O: "a ship" = "pire"? No, "a boat" = "pire" in R1 \#1?} \\
\texttt{                      {-} If OSV: O="a ship", S="a boat", V="watch".} \\
\texttt{                      {-} Translation: "pire deniru dugu".} \\
\texttt{                      {-} If O="pire", S="dugu", V="deniru".} \\
\texttt{                      {-} Then "a ship" = "pire", "a boat" = "dugu".} \\
\texttt{                    {-} R1 \#2: "a boat watch two ships" => "pirenu nu deniru dugu".} \\
\texttt{                      {-} O="two ships", S="a boat", V="watch".} \\
\texttt{                      {-} Translation: "pirenu nu deniru dugu".} \\
\texttt{                      {-} If O="pirenu nu", S="dugu", V="deniru".} \\
\texttt{                      {-} Then "two ships" = "pirenu nu". "a boat" = "dugu".} \\
\texttt{                      {-} Check R1 \#1: "a ship" = "pire". "a boat" = "dugu".} \\
\texttt{                      {-} Check R1 \#3: "two boats watch a ship" => "pire deniru dugunu nu".} \\
\texttt{                        {-} O="a ship", S="two boats", V="watch".} \\
\texttt{                        {-} Translation: "pire deniru dugunu nu".} \\
\texttt{                        {-} If O="pire", S="dugunu nu", V="deniru".} \\
\texttt{                        {-} Then "a ship" = "pire". "two boats" = "dugunu nu".} \\
\texttt{                        {-} Consistency check:} \\
\texttt{                          {-} "a ship" = "pire" (R1 \#1, R1 \#3).} \\
\texttt{                          {-} "a boat" = "dugu" (R1 \#1, R1 \#2).} \\
\texttt{                          {-} "two ships" = "pirenu nu" (R1 \#2).} \\
\texttt{                          {-} "two boats" = "dugunu nu" (R1 \#3, R1 \#4).} \\
\texttt{                          {-} "a dog" = ?} \\
\texttt{                          {-} "three dogs" = "tipi puga tipinu ga" (R2 \#2).} \\
\texttt{                            {-} If OSV: O="a dog", S="three dogs", V="hear".} \\
\texttt{                            {-} Translation: "tipi puga tipinu ga".} \\
\texttt{                            {-} If O="tipi puga tipinu ga"? Too long.} \\
\texttt{                            {-} Maybe "hear" is "ga"?} \\
\texttt{                            {-} "three dogs" = "tipi puga tipinu"? "a dog" = "ga"?} \\
\texttt{                            {-} Let's check R2 \#3: "three dogs carry three ships" => "pirenu ga geti tipinu ga".} \\
\texttt{                              {-} O="three ships", S="three dogs", V="carry".} \\
\texttt{                              {-} Translation: "pirenu ga geti tipinu ga".} \\
\texttt{                              {-} If OSV: O="pirenu ga", S="tipinu ga", V="geti".} \\
\texttt{                              {-} "three ships" = "pirenu ga".} \\
\texttt{                              {-} "three dogs" = "tipinu ga".} \\
\texttt{                              {-} "carry" = "geti".} \\
\texttt{                              {-} Check R2 \#2: "three dogs hear a dog" => "tipi puga tipinu ga".} \\
\texttt{                                {-} O="a dog", S="three dogs", V="hear".} \\
\texttt{                                {-} Translation: "tipi puga tipinu ga".} \\
\texttt{                                {-} If OSV: O="tipi puga tipinu ga"? No.} \\
\texttt{                                {-} Maybe "hear" = "puga"?} \\
\texttt{                                {-} "three dogs" = "tipi puga tipinu"? "a dog" = "ga"?} \\
\texttt{                                {-} "tipi puga tipinu ga" {-}> O="tipi puga tipinu", S="ga"? No.} \\
\texttt{                                {-} Maybe the word order is **Subject{-}Object{-}Verb (SOV)**?} \\
\texttt{                                  {-}} \\
            }
        }
    }
    & & \\ \\

    \theutterance \stepcounter{utterance}  
    & & & \multicolumn{2}{p{0.3\linewidth}}{
        \cellcolor[rgb]{0.9,0.9,0.9}{
            \makecell[{{p{\linewidth}}}]{
                \texttt{\tiny{[GM$|$GM]}}
                \texttt{could not find any probe sentences; put them under a 'PROBES:' line, one per line} \\
            }
        }
    }
    & & \\ \\

    \theutterance \stepcounter{utterance}  
    & & & \multicolumn{4}{p{0.6\linewidth}}{
        \cellcolor[rgb]{0.9,0.9,0.9}{
            \makecell[{{p{\linewidth}}}]{
                \texttt{\tiny{[P1$\langle$GM]}}
                \texttt{Your last response could not be processed: could not find any probe sentences; put them under a 'PROBES:' line, one per line.} \\
\texttt{List your probes (a brief THOUGHTS line is optional):} \\
\texttt{PROBES:} \\
\texttt{1. <an English sentence>} \\
\texttt{(up to 3 probe lines; at least one)} \\
            }
        }
    }
     \\ \\

    \theutterance \stepcounter{utterance}  
    & \multicolumn{4}{p{0.6\linewidth}}{
        \cellcolor[rgb]{0.561,0.737,0.561}{
            \makecell[{{p{\linewidth}}}]{
                \texttt{\tiny{[P1$\rangle$GM]}}
                \texttt{THOUGHTS: Based on the previous translations, I have deduced the following:} \\
\texttt{1.  **Word Order**: The language is **Object{-}Subject{-}Verb (OSV)**.} \\
\texttt{    *   Evidence: "a boat watch a ship" => "pire deniru dugu".} \\
\texttt{        *   If OSV: Object="a ship" (pire), Subject="a boat" (dugu), Verb="watch" (deniru).} \\
\texttt{        *   Check with "two dogs chase two boats" => "dugunu nu buro tipinu nu".} \\
\texttt{        *   If OSV: Object="two boats" (dugunu nu), Subject="two dogs" (tipinu nu), Verb="chase" (buro).} \\
\texttt{        *   This fits perfectly. "dugunu nu" is the object (two boats), "tipinu nu" is the subject (two dogs).} \\
\texttt{2.  **Noun Meanings \& Plurals**:} \\
\texttt{    *   "a ship" = `pire` (singular). "two ships" = `pirenu nu` (plural marker `{-}nu` added to noun, then space + `nu`? Or `pire` + `nu` + `nu`? Let's re{-}examine).} \\
\texttt{    *   Wait, let's re{-}parse "a boat watch a ship" => `pire deniru dugu`.} \\
\texttt{        *   If OSV: O=`pire` (ship), S=`dugu` (boat).} \\
\texttt{        *   So `pire` = ship, `dugu` = boat.} \\
\texttt{    *   "a boat watch two ships" => `pirenu nu deniru dugu`.} \\
\texttt{        *   O=`two ships` = `pirenu nu`. S=`a boat` = `dugu`.} \\
\texttt{        *   Pattern for "two X": `X` + `nu` + `nu`? Or `X` changes to `Xnu` then `nu`?} \\
\texttt{        *   "two dogs chase two boats" => `dugunu nu buro tipinu nu`.} \\
\texttt{        *   O=`two boats` = `dugunu nu`. S=`two dogs` = `tipinu nu`.} \\
\texttt{        *   "a dog" in "three dogs hear a dog" => `tipi puga tipinu ga`.} \\
\texttt{        *   If OSV: O=`a dog` = `tipi puga tipinu ga`? No, that's the whole sentence.} \\
\texttt{        *   Let's look at "three dogs hear a dog" => `tipi puga tipinu ga`.} \\
\texttt{        *   If OSV: O=`a dog`, S=`three dogs`, V=`hear`.} \\
\texttt{        *   Translation: `tipi puga tipinu ga`.} \\
\texttt{        *   Maybe `tipi` = three, `puga` = dogs? No, "three dogs" is the subject.} \\
\texttt{        *   Let's try: O=`a dog` = `tipi`. S=`three dogs` = `puga tipinu ga`? No.} \\
\texttt{        *   Alternative: Maybe `ga` is the verb? No, `deniru` is watch.} \\
\texttt{        *   Let's reconsider the OSV hypothesis with "three dogs hear a dog".} \\
\texttt{        *   If O=`a dog`, S=`three dogs`, V=`hear`.} \\
\texttt{        *   Translation: `tipi puga tipinu ga`.} \\
\texttt{        *   Could it be `tipi` (a dog) `puga` (hear) `tipinu ga` (three dogs)?} \\
\texttt{        *   If so, O=`tipi`, V=`puga`, S=`tipinu ga`.} \\
\texttt{        *   Then "three dogs" = `tipinu ga`. "a dog" = `tipi`.} \\
\texttt{        *   Check "three dogs carry three ships" => `pirenu ga geti tipinu ga`.} \\
\texttt{        *   If OSV: O=`three ships` = `pirenu ga`. S=`three dogs` = `tipinu ga`. V=`carry` = `geti`.} \\
\texttt{        *   This works! `pirenu ga` = three ships. `tipinu ga` = three dogs.} \\
\texttt{        *   So `ga` seems to be a plural marker for definite or specific groups? Or maybe `ga` is part of the number phrase?} \\
\texttt{        *   "three" = `tipi`? No, "three dogs" = `tipinu ga`. "three ships" = `pirenu ga`.} \\
\texttt{        *   Maybe `tipi` is the root for "three"? `tipi` + `nu` + `ga`?} \\
\texttt{        *   Let's look at "two dogs" = `tipinu nu`. "two boats" = `dugunu nu`.} \\
\texttt{        *   "two" = `dugu`? No, "a boat" = `dugu`.} \\
\texttt{        *   "two" = `dugunu`? No.} \\
\texttt{        *   Let's look at the number words again.} \\
\texttt{        *   "two dogs" = `tipinu nu`. "three dogs" = `tipinu ga`.} \\
\texttt{        *   Difference is `nu` vs `ga`.} \\
\texttt{        *   "two boats" = `dugunu nu`. "three ships" = `pirenu ga`.} \\
\texttt{        *   "two ships" = `pirenu nu`.} \\
\texttt{        *   It seems `nu` marks "two" and `ga` marks "three"?} \\
\texttt{        *   But "three dogs" = `tipinu ga`. "three ships" = `pirenu ga`.} \\
\texttt{        *   "two dogs" = `tipinu nu`. "two boats" = `dugunu nu`. "two ships" = `pirenu nu`.} \\
\texttt{        *   So `nu` = two. `ga` = three?} \\
\texttt{        *   Wait, "three farmers watch three farmers" => `nidinu ga deniru nidinu ga`.} \\
\texttt{        *   "three farmers" = `nidinu ga`.} \\
\texttt{        *   "two roads hold three trees" => `gutinu ga nira dobudinu nu`.} \\
\texttt{        *   O=`three trees` = `dobudinu nu`. S=`two roads` = `gutinu ga`.} \\
\texttt{        *   If `nu` = two, then `dobudinu nu` = "two trees"? But English says "three trees".} \\
\texttt{        *   Contradiction.} \\
\texttt{        *   Let's re{-}evaluate "two roads hold three trees" => `gutinu ga nira dobudinu nu`.} \\
\texttt{        *   If OSV: O=`three trees` = `dobudinu nu`. S=`two roads` = `gutinu ga`.} \\
\texttt{        *   If `nu` = two, then `dobudinu nu` means "two trees". But English is "three trees".} \\
\texttt{        *   Maybe `nu` is the plural marker for indefinite?} \\
\texttt{        *   "two X" = `X` + `nu` + `nu`?} \\
\texttt{        *   "two dogs" = `tipinu nu`. "tipi" = dog? "nu" = two?} \\
\texttt{        *   "three dogs" = `tipinu ga`. "ga" = three?} \\
\texttt{        *   "two ships" = `pirenu nu`. "pire" = ship.} \\
\texttt{        *   "three ships" = `pirenu ga`.} \\
\texttt{        *   "two boats" = `dugunu nu`. "dugu" = boat.} \\
\texttt{        *   "three trees" = `dobudinu nu`. "dobudi" = tree?} \\
\texttt{        *   "two roads" = `gutinu ga`. "guti" = road?} \\
\texttt{        *   This implies `ga` = three, `nu` = two.} \\
\texttt{        *   But "three trees" = `dobudinu nu`. Why `nu`?} \\
\texttt{        *   Maybe the number word comes *after* the noun?} \\
\texttt{        *   "two dogs" = `tipi` (dog) `nu` (two)? No, `tipinu nu`.} \\
\texttt{        *   Maybe `tipi` = three? `tipinu` = three dogs? `ga` = ?} \\
\texttt{        *   Let's look at "three dogs hear a dog" => `tipi puga tipinu ga`.} \\
\texttt{        *   If OSV: O=`a dog` = `tipi`. V=`puga` (hear). S=`three dogs` = `tipinu ga`.} \\
\texttt{        *   "three dogs" = `tipinu ga`.} \\
\texttt{        *   "three ships" = `pirenu ga`.} \\
\texttt{        *   "three farmers" = `nidinu ga`.} \\
\texttt{        *   "three trees" = `dobudinu nu`. (Wait, why `nu` here?)} \\
\texttt{        *   "two roads" = `gutinu ga`.} \\
\texttt{        *   "two dogs" = `tipinu nu`.} \\
\texttt{        *   "two boats" = `dugunu nu`.} \\
\texttt{        *   "two ships" = `pirenu nu`.} \\
\texttt{        *   Pattern:} \\
\texttt{            *   Two: `...nu` (at the end of the noun phrase).} \\
\texttt{            *   Three: `...ga` (at the end of the noun phrase).} \\
\texttt{        *   But "three trees" = `dobudinu nu`. This breaks the pattern.} \\
\texttt{        *   Unless "three trees" is actually "two trees"? No, English says "three".} \\
\texttt{        *   Maybe `dobudinu` is the word for "three trees"?} \\
\texttt{        *   Let's check "two roads hold three trees".} \\
\texttt{        *   S=`two roads` = `gutinu ga`. (Roads = `guti`? Two = `nu`? `ga` = ?)} \\
\texttt{        *   O=`three trees` = `dobudinu nu`. (Trees = `dobudi`? Three = `nu`? Two = `ga`?)} \\
\texttt{        *   If `ga` = two and `nu` = three:} \\
\texttt{            *   "two roads" = `gutinu ga` {-}> `guti` + `nu` + `ga`? No.} \\
\texttt{            *   "three trees" = `dobudinu nu` {-}> `dobudi` + `nu` + `nu`?} \\
\texttt{        *   Let's try: Number words are suffixes.} \\
\texttt{        *   "two" = `nu`. "three" = `ga`.} \\
\texttt{        *   "two dogs" = `dog` + `nu` + `nu`? `tipinu nu`.} \\
\texttt{        *   "three dogs" = `dog` + `nu` + `ga`? `tipinu ga`.} \\
\texttt{        *   "two ships" = `ship` + `nu` + `nu`? `pirenu nu`.} \\
\texttt{        *   "three ships" = `ship` + `nu` + `ga`? `pirenu ga`.} \\
\texttt{        *   "two boats" = `boat` + `nu` + `nu`? `dugunu nu`.} \\
\texttt{        *   "three trees" = `tree` + `nu` + `ga`? `dobudinu ga`.} \\
\texttt{        *   But the translation is `dobudinu nu`.} \\
\texttt{        *   Is it possible "three trees" is a typo in my analysis?} \\
\texttt{        *   "two roads hold three trees" => `gutinu ga nira dobudinu nu`.} \\
\texttt{        *   Maybe `dobudinu` means "three trees" and `nu` is the plural marker?} \\
\texttt{        *   Or maybe `nu` is "two" and `ga` is "three", but "three trees" uses `nu` because of agreement?} \\
\texttt{        *   Let's look at "three farmers watch three farmers" => `nidinu ga deniru nidinu ga`.} \\
\texttt{        *   "three farmers" = `nidinu ga`.} \\
\texttt{        *   "two roads" = `gutinu ga`.} \\
\texttt{        *   "two dogs" = `tipinu nu`.} \\
\texttt{        *   "two boats" = `dugunu nu`.} \\
\texttt{        *   "two ships" = `pirenu nu`.} \\
\texttt{        *   "three trees" = `dobudinu nu`.} \\
\texttt{        *   Hypothesis:} \\
\texttt{            *   `nu` = two.} \\
\texttt{            *   `ga` = three.} \\
\texttt{            *   BUT "three trees" = `dobudinu nu`. This is the outlier.} \\
\texttt{            *   Maybe `dobudinu` is the word for "three trees" and `nu` is the plural marker for indefinite?} \\
\texttt{            *   "two roads" = `gutinu ga`. `gutinu` = two roads? `ga` = ?} \\
\texttt{            *   "two dogs" = `tipinu nu`. `tipinu` = two dogs? `nu` = ?} \\
\texttt{            *   Maybe the number word is the *last* word of the NP.} \\
\texttt{            *   "two dogs" = `tipinu` (two dogs) `nu` (plural marker?).} \\
\texttt{            *   "three dogs" = `tipinu` (three dogs) `ga` (plural marker?).} \\
\texttt{            *   "three trees" = `dobudinu` (three trees) `nu` (plural marker?).} \\
\texttt{            *   "two roads" = `gutinu` (two roads) `ga` (plural marker?).} \\
\texttt{            *   This implies `nu` and `ga` are plural markers, not numbers.} \\
\texttt{            *   Then where are the numbers?} \\
\texttt{            *   "two dogs" vs "three dogs".} \\
\texttt{            *   `tipinu nu` vs `tipinu ga`.} \\
\texttt{            *   Maybe `nu` = two, `ga` = three.} \\
\texttt{            *   Then "three trees" = `dobudinu nu` means "two trees"?} \\
\texttt{            *   But the English is "three trees".} \\
\texttt{            *   Is it possible the informant made a mistake? No.} \\
\texttt{            *   Is it possible `dobudinu` means "three trees" and `nu` is just a plural marker that appears on all plurals?} \\
\texttt{            *   If so, "two roads" = `gutinu ga`. `ga` would be the plural marker?} \\
\texttt{            *   Then `nu` = two, `ga` = three?} \\
\texttt{            *   "two dogs" = `tipinu nu`. `nu` = two.} \\
\texttt{            *   "three dogs" = `tipinu ga`. `ga` = three.} \\
\texttt{            *   "two ships" = `pirenu nu`. `nu` = two.} \\
\texttt{            *   "three ships" = `pirenu ga`. `ga` = three.} \\
\texttt{            *   "two boats" = `dugunu nu`. `nu` = two.} \\
\texttt{            *   "three trees" = `dobudinu nu`. `nu` = two? But English says three.} \\
\texttt{            *   Maybe `dobudinu` is "three trees" and `nu` is the plural marker?} \\
\texttt{            *   And `gutinu` is "two roads" and `ga` is the plural marker?} \\
\texttt{            *   This would mean `nu` and `ga` are plural markers, and the number is encoded in the root?} \\
\texttt{            *   "tipi" = dog. "tipinu" = two dogs? "tipinu ga" = three dogs?} \\
\texttt{            *   "pire" = ship. "pirenu" = two ships? "pirenu ga" = three ships?} \\
\texttt{            *   "dugu" = boat. "dugunu" = two boats?} \\
\texttt{            *   "dobudi" = tree. "dobudinu" = three trees?} \\
\texttt{            *   "guti" = road. "gutinu" = two roads?} \\
\texttt{            *   "nidi" = farmer. "nidinu" = three farmers?} \\
\texttt{            *   This seems inconsistent.} \\
\texttt{            *   Let's try: Number words are `nu` (two) and `ga` (three).} \\
\texttt{            *   "two X" = `X` + `nu` + `nu`?} \\
\texttt{            *   "three X" = `X` + `nu` + `ga`?} \\
\texttt{            *   "two dogs" = `tipi` + `nu` + `nu` = `tipinu nu`.} \\
\texttt{            *   "three dogs" = `tipi` + `nu` + `ga` = `tipinu ga`.} \\
\texttt{            *   "two ships" = `pire` + `nu` + `nu` = `pirenu nu`.} \\
\texttt{            *   "three ships" = `pire` + `nu` + `ga` = `pirenu ga`.} \\
\texttt{            *   "two boats" = `dugu` + `nu` + `nu` = `dugunu nu`.} \\
\texttt{            *   "three trees" = `dobudi` + `nu` + `ga` = `dobudinu ga`.} \\
\texttt{            *   But the translation is `dobudinu nu`.} \\
\texttt{            *   Is it possible "three trees" is actually "two trees" in the informant's mind? No.} \\
\texttt{            *   Is it possible `dobudinu` means "three trees" and `nu` is the plural marker?} \\
\texttt{            *   And `gutinu` means "two roads" and `ga` is the plural marker?} \\
\texttt{            *   Then `nu` and `ga` are just plural markers.} \\
\texttt{            *   But then how do we distinguish "two" and "three"?} \\
\texttt{            *   Maybe the root changes?} \\
\texttt{            *   "tipi" = dog. "tipinu" = two dogs. "tipinu ga" = three dogs.} \\
\texttt{            *   "pire" = ship. "pirenu" = two ships. "pirenu ga" = three ships.} \\
\texttt{            *   "dugu" = boat. "dugunu" = two boats.} \\
\texttt{            *   "dobudi" = tree. "dobudinu" = three trees.} \\
\texttt{            *   "guti" = road. "gutinu" = two roads.} \\
\texttt{            *   "nidi" = farmer. "nidinu" = three farmers.} \\
\texttt{            *   This implies:} \\
\texttt{                *   `nu` = two.} \\
\texttt{                *   `ga` = three.} \\
\texttt{                *   BUT "three trees" = `dobudinu nu`. This has `nu` at the end.} \\
\texttt{                *   Maybe `dobudinu` is "three trees" and `nu` is the plural marker?} \\
\texttt{                *   And `gutinu` is "two roads" and `ga` is the plural marker?} \\
\texttt{                *   This is contradictory.} \\
\texttt{            *   Let's reconsider the "three trees" translation: `dobudinu nu`.} \\
\texttt{            *   Maybe `dobudinu` = "three trees" and `nu` = plural marker?} \\
\texttt{            *   And `gutinu` = "two roads" and `ga` = plural marker?} \\
\texttt{            *   Then `nu` and `ga` are plural markers.} \\
\texttt{            *   But then how do we know "two" vs "three"?} \\
\texttt{            *   Maybe the root `dobudi` implies "three"? No.} \\
\texttt{            *   Maybe `dobudinu` is "two trees"? No, English says "three".} \\
\texttt{            *   Is it possible `nu` = two, `ga` = three, and "three trees" is a mistake in my parsing?} \\
\texttt{            *   "two roads hold three trees" => `gutinu ga nira dobudinu nu`.} \\
\texttt{            *   S=`two roads` = `gutinu ga`.} \\
\texttt{            *   O=`three trees` = `dobudinu nu`.} \\
\texttt{            *   If `ga` = three, then `gutinu ga` = "two roads" (with `ga` as plural marker?).} \\
\texttt{            *   If `nu` = two, then `dobudinu nu` = "three trees" (with `nu` as plural marker?).} \\
\texttt{            *   This suggests `ga` and `nu` are plural markers, and the number is in the root?} \\
\texttt{            *   "tipi" = dog. "tipinu" = two dogs. "tipinu ga" = three dogs.} \\
\texttt{            *   "pire" = ship. "pirenu" = two ships. "pirenu ga" = three ships.} \\
\texttt{            *   "dugu" = boat. "dugunu" = two boats.} \\
\texttt{            *   "dobudi" = tree. "dobudinu" = three trees.} \\
\texttt{            *   "guti" = road. "gutinu" = two roads.} \\
\texttt{            *   "nidi" = farmer. "nidinu" = three farmers.} \\
\texttt{            *   This fits if:} \\
\texttt{                *   `nu` = two.} \\
\texttt{                *   `ga` = three.} \\
\texttt{                *   BUT "three trees" = `dobudinu nu`. This has `nu` at the end.} \\
\texttt{                *   Maybe `dobudinu` is "three trees" and `nu` is the plural marker?} \\
\texttt{                *   And `gutinu` is "two roads" and `ga` is the plural marker?} \\
\texttt{                *   This is still contradictory.} \\
\texttt{            *   Let's try: `nu` = two, `ga` = three.} \\
\texttt{            *   "two dogs" = `tipi` + `nu` + `nu` = `tipinu nu`.} \\
\texttt{            *   "three dogs" = `tipi` + `nu` + `ga` = `tipinu ga`.} \\
\texttt{            *   "two ships" = `pire` + `nu` + `nu` = `pirenu nu`.} \\
\texttt{            *   "three ships" = `pire` + `nu` + `ga` = `pirenu ga`.} \\
\texttt{            *   "two boats" = `dugu` + `nu` + `nu` = `dugunu nu`.} \\
\texttt{            *   "three trees" = `dobudi` + `nu` + `ga` = `dobudinu ga`.} \\
\texttt{            *   But the translation is `dobudinu nu`.} \\
\texttt{            *   Is it possible `dobudinu` means "three trees" and `nu` is the plural marker?} \\
\texttt{            *   And `gutinu` means "two roads" and `ga` is the plural marker?} \\
\texttt{            *   Then `nu` and `ga` are plural markers.} \\
\texttt{            *   But then how do we distinguish "two" and "three"?} \\
\texttt{            *   Maybe the root changes?} \\
\texttt{            *   "tipi" = dog. "tipinu" = two dogs. "tipinu ga" = three dogs.} \\
\texttt{            *   "pire" = ship. "pirenu" = two ships. "pirenu ga" = three ships.} \\
\texttt{            *   "dugu" = boat. "dugunu" = two boats.} \\
\texttt{            *   "dobudi" = tree. "dobudinu" = three trees.} \\
\texttt{            *   "guti" = road. "gutinu" = two roads.} \\
\texttt{            *   "nidi" = farmer. "nidinu" = three farmers.} \\
\texttt{            *   This fits if:} \\
\texttt{                *   `nu` = two.} \\
\texttt{                *   `ga` = three.} \\
            }
        }
    }
    & & \\ \\

    \theutterance \stepcounter{utterance}  
    & & & \multicolumn{2}{p{0.3\linewidth}}{
        \cellcolor[rgb]{0.9,0.9,0.9}{
            \makecell[{{p{\linewidth}}}]{
                \texttt{\tiny{[GM$|$GM]}}
                \texttt{could not find any probe sentences; put them under a 'PROBES:' line, one per line} \\
            }
        }
    }
    & & \\ \\

\end{supertabular}
}

\end{document}
